\subsection{Discussion and conclusions}
\label{sec:catalog:conclusions}
%%%%


We have presented an approach to statistically push the \Fermi-LAT sensitivity to point-like sources at high latitudes below  the current threshold for detection, leveraging on the fact that the underlying source count  distribution function is constrained even for lower test statistics ($TS$). This function has been recently re-derived using machine learning methods, yielding results in  accordance with existing pixel statistics analyses. This information, not accounted for in the current catalogue construction procedure, allowed us to provide a catalogue of directions (i.e.~firing pixels) in the sky likely to be associated to sources, albeit only in a probabilistic sense, that we assessed via a Kolmogorov-Smirnov (KS) test. The actual catalogue of directions depends on meta-parameters such as the confidence level $1-\alpha$ of the KS test, or the fraction $QF$ of simulated skies for which the simulation agrees with data above a certain $TS$ level. For reasonable choices of these meta-parameters, we found a number of ``firing pixels'' $\sim$50\%  higher than what one would infer from a catalogue only including sources listed in the latest incarnation of the \Fermi-LAT catalogue. Our results also pass some sanity checks, like the fact that \Fermi-LAT catalogue sources which are luminous enough are found within the directions inferred by our method, within the angular resolution of the instrument. If compared with the ``local'' iterative procedure followed by \Fermi-LAT to establish their catalogue, another possible advantage of our catalogue of directions in the sky is that its $TS$ scale has a homogeneous meaning, and may be more suitable for global statistical analyses.  

We provide our results in digital form, allowing the user to select the meta-parameters $QF$ and $\alpha$ to be more or less restrictive in their selection criteria, hopefully adapting to a wide range of applications. The most obvious one we can think of consists of cross-correlation studies, either multi-wavelength or multi-messenger, where the statistical advantage of  a significantly larger sample more than compensates having a (sub-leading) fraction of spurious directions. 
Another straightforward use of our results could be a guided source search and identification program via multi-wavelength studies, which would also help transforming some of these candidate source directions into  \textit{bona fide} sources. In turn, these studies may further benefit of machine learning techniques that have been recently proposed to ease the  probabilistic classification of unassociated sources~\cite{Bhat:2022}.

We stress that, although our results are based on the source-count distribution as inferred from a specific machine learning analysis of the photon counts, the method developed here to identify ``firing pixels'' can be applied to any other determination of the $\dNdS$ of the high-latitude sky, and it is, in this respect, independent on the choice of the $\dNdS$ as long as the latter correctly describes point sources in the faint regime.

Since our work had primarily a methodological proof-of-principle motivation, we have resorted to a couple of technical simplifications: We worked with an energy-integrated spectrum and we limited ourselves to high Galactic latitudes. We could rather naturally lift the former approximation  by a brute force approach. We expect the main complication to be computational, but the problem should remain affordable in a reasonable time. Also the latter limitation could be in principle lifted, extending the analysis to low Galactic latitudes. The main problem we can anticipate is that the existing background models fail to perform as well at low latitudes as they do at high latitudes. One also needs to include multiple populations contributing to both resolved and unresolved emission, see e.g.~\cite{Calore:2014oga} and sec. 6.2 in ~\cite{Fermi-LAT:2015bhf} for the role played by Galactic millisecond pulsars in this context. In the worst case, the method may not allow one to gain much more with respect to more traditional techniques. Depending on the outcome, one may however think of multi-zone models for the background or the like.  

A third direction for future progress is to attach a  more realistic probability scale to the $TS$ maps. We have constructed our synthetic skies only based on the statistical properties for the sources previously deduced via the pixel analysis. However, only a subset of these simulations is consistent with the properties of the actual sky: Notably at high flux where only a few sources exist, the associated $TS$ distribution significantly depends on the distribution of the sources in the sky randomly falling in a lower or higher background region (for an illustration, see the large variability at high $TS$ in figure~\ref{fig:1}). As briefly discussed in the article, a promising approach would be to narrow the sample of admissible synthetic skies used to compute $QF$ to the sole simulations which are statistically consistent with the $TS$ pattern of the most luminous sources. Such `constrained' simulations should then improve the reliability of the $QF$ scale. Another working direction with the same goal could be to improve the description of background and/or the parameterization for sub-threshold sources in the pixel analyses beyond what considered in~\cite{Amerio:2023uet}. This may bring the simulation distribution illustrated in figure~\ref{fig:1} in even better agreement with the $TS$ histogram. Finally, exploring alternative  metrics to the agreement in the KS sense is also a direction potentially worth exploring. 

%\medskip

